\clearpage
\VerdierPhysicalPageAnchor{0180}%
\begingroup\small\setlength{\parskip}{0pt}\setlength{\emergencystretch}{1.5em}
\setlength{\abovedisplayskip}{4pt plus 2pt minus 2pt}%
\setlength{\belowdisplayskip}{4pt plus 2pt minus 2pt}%

\VerdierSemanticAnchor{ch3d-1-2-8-duality-cont}
Henceforth, unless otherwise stated, we identify the categories
$D(\mathcal A)^o$ and $D(\mathcal A^o)$ by means of the isomorphism $i$.
Note that this identification identifies $D^b(\mathcal A)^o$ with
$D^b(\mathcal A^o)$, $D^+(\mathcal A)^o$ with $D^-(\mathcal A^o)$, and
$D^-(\mathcal A)^o$ with $D^+(\mathcal A^o)$.

\VerdierSemanticAnchor{ch3d-1-2-9}
\medskip\noindent\textbf{1.2.9.}\quad
The canonical functor $\mathcal A\longrightarrow K(\mathcal A)$
(Chap. II, 1.3.3), composed with the functor
$Q\colon K(\mathcal A)\longrightarrow D(\mathcal A)$, gives a fully
faithful functor $\mathcal A\longrightarrow D(\mathcal A)$: indeed, let
$X$ and $Y$ be two objects of $\mathcal A$. A morphism from $X$ to $Y$ in
$D(\mathcal A)$ is (Chap. II, 2.2.2) the class $\overline D$ of a diagram in
$K(\mathcal A)$:

\VerdierSemanticAnchor{ch3d-1-2-9-roof}
\[
\overline D=X\xleftarrow{\ s\ }Z^\bullet\xrightarrow{\ m\ }Y,
\]
where $s$ is a quasi-isomorphism in $K(\mathcal A)$. Consequently,
$s\colon Z^\bullet\longrightarrow X$ induces an isomorphism
$H^0(s)\colon H^0(Z^\bullet)\xrightarrow{\sim}H(X)=X$, and one immediately
checks that $\overline D=Q\bigl(H^0(m)H^0(s)^{-1}\bigr)$. This proves that
the homomorphism
\[
\operatorname{Hom}_{\mathcal A}(X,Y)\longrightarrow
\operatorname{Hom}_{D(\mathcal A)}(X,Y)
\]
is surjective. Let us prove that it is injective. A morphism
$f\colon X\longrightarrow Y$ in $\mathcal A$ defines the zero morphism in
$D(\mathcal A)$ if and only if there exists a quasi-isomorphism
$s\colon Z^\bullet\longrightarrow X$ in $K(\mathcal A)$ such that $fs=0$.
It follows that $H^0(f)H^0(s)=0$. Since $H^0(s)$ is an isomorphism, it
follows that $H^0(f)=f=0$.

Note also that the functor $\mathcal A\longrightarrow D(\mathcal A)$ is
injective on the sets of objects and therefore realizes $\mathcal A$ as a
full subcategory of $D(\mathcal A)$. Henceforth, the category $\mathcal A$
will always be regarded as a subcategory of its derived category by means of
this functor. The canonical cohomological functor
$D(\mathcal A)\xrightarrow{H}\mathcal A$, restricted to the subcategory
$\mathcal A$, is the identity functor.

\VerdierSemanticAnchor{ch3d-1-2-10}
\medskip\noindent\textbf{Proposition 1.2.10.}\quad
\textit{a) An object $X$ of $D(\mathcal A)$ is isomorphic to an object of
$\mathcal A$ if and only if $H^i(X)=0$ for every $i\neq0$. It is then
canonically isomorphic to $H^0(X)$.}

\noindent\textit{b) Let $X$ and $Y$ be two objects of $\mathcal A$. Then}
\[
\operatorname{Ext}_{\mathcal A}^{n}(X,Y)=0\quad\text{for }n<0,
\qquad
\operatorname{Ext}_{\mathcal A}^{0}(X,Y)
=\operatorname{Hom}_{\mathcal A}(X,Y).
\]

\VerdierSemanticAnchor{ch3d-1-2-10-proof}
Let us prove a). It is clear that $H^i(X)=0$ for $i\neq0$ when $X$ is an
object of $\mathcal A$. Let $X$ be an object of $D(\mathcal A)$ such that
$H^i(X)$ vanishes for $i\neq0$.
\endgroup
